%==============================================================
%  Chapter 4 -- Embedding-Based Item-to-Item Recommendation
%  Sources: "Hierarchical Cosine-Softmax Objective for Product
%  Embeddings" (spec) and "hierarchical_embedding_experiments"
%  (as-built + experiments). Both are the user's own work.
%
%  NOTE: all reported numbers are TRAINING-OBJECTIVE results
%  (weighted NLL). Recommendation-quality (offline CTR) evaluation
%  is out of scope for this chapter and stated as such.
%==============================================================
\chapter{Embedding-Based Item-to-Item Recommendation}
\label{ch:embedding}

The chapters so far model recommendation through transition matrices and their
powers. This chapter takes a different representation: it learns a single dense
\emph{embedding} for every product, such that products between which users
transition end up close together in vector space, and then serves
recommendations by nearest-neighbour search in that space. The construction is
at once content-based and collaborative --- embeddings are \emph{cold-started}
from product text metadata and then \emph{refined} from observed item-to-item
interactions --- and it is organised around the product category hierarchy. We
present the model as originally specified (a hierarchical cosine-softmax
objective), the extensions made during development (trainable centroids and a
flat alternative), the implementation as actually built, and a sequence of
controlled training experiments together with what they reveal.

A note on scope frames the chapter. The experiments of
Section~\ref{sec:emb-experiments} and the findings drawn from them
(Section~\ref{sec:emb-findings}) are \emph{training-objective} results --- the
weighted negative log-likelihood being optimised --- and their conclusions are
statements about optimisation behaviour. Whether a lower training loss actually
yields better recommendations is a separate question, taken up in
Section~\ref{sec:emb-ctr}, which reports offline click-through evaluation. As we
shall see, the two do not always agree, and in two instances they point in
opposite directions.

%--------------------------------------------------------------
\section{Problem and Setup}
\label{sec:emb-problem}

The goal is an item-to-item recommender for IndiaMART that learns a dense
embedding $e_p \in \mathbb{R}^{d}$ for every product $p$, using observed
interaction data together with the product category hierarchy. Two data
structures are available:

\begin{itemize}
  \item \textbf{Interaction logs} of the form $(\text{prod } i \to
        \text{prod } j)$, recording that users transitioned from product $i$ to
        product $j$.
  \item \textbf{A hierarchical taxonomy},
        $\text{Industry} \to \text{Sub-category} \to \text{Main-category} \to
        \text{Product}$. The core formulation uses the two levels immediately
        above products --- sub-category and main-category.
  \item \textbf{Product-level metadata} (name, summary, category text,
        specifications, seller details) used \emph{only} to initialise product
        embeddings.
\end{itemize}

Training is confined to a clean slice of the logs. Only interactions with
$\text{MODID}=\text{IMOB}$ are used; only rows with positive $\text{ProdID}_1$
and $\text{ProdID}_2$ are considered valid; an interaction is dropped if either
product lacks an initial embedding; the taxonomy carried by the interaction
data is treated as the source of truth during training; and duplicate
transitions are aggregated, so training operates on \emph{weighted} interaction
pairs. The objectives are then to (i)~learn one product embedding $e_p$ per
product, (ii)~use a hierarchical, cosine-similarity-based softmax factorisation
over sub-category, main-category, and products to define a likelihood for each
observed transition $(i\to j)$, (iii)~maximise the log-likelihood
(equivalently, minimise the negative log-likelihood) by gradient descent, and
(iv)~in the baseline, use fixed frozen centroids for sub- and main-categories,
computed from the initial product embeddings.

At inference the learned embeddings are consumed by a vector-search setup that
computes the top-$k$ similar products for each product. Both the products
updated through interaction training and those that remain at their initial
text-based embeddings are retained in the final retrieval space. The hierarchy
is thus used during \emph{training} as a structured objective, not for online
inference.

\subsection{Content-Based Product Vectors}
\label{sec:emb-content-init}
The initial embedding of every product is built from its catalogue metadata in
the product description dataset (Section~\ref{sec:product-desc}). For each
catalogue entry a single text document is assembled by concatenating the product
name, its master category, general category, highest-level product group, and
primary product category, together with the product specifications. The
specifications, held as JSON in the \texttt{custom\_specs\_json} field, are
flattened into text as a list of attribute--value pairs, which keeps their
content while making them readable by a text encoder. The result is one document
per product that combines the product's identity, its place in the category
hierarchy, and its detailed specifications.

Each document is encoded with the \texttt{all-MiniLM-L6-v2} sentence-transformer,
which maps a variable-length text to a $384$-dimensional dense vector in a shared
semantic space; this is where $d=384$ throughout the chapter comes from. Products
with similar descriptions are placed near one another, so these vectors carry
content similarity even when interaction data is sparse or absent --- the
cold-start signal the model starts from. Because the catalogue metadata does not
change during training, the vectors are computed once during preprocessing and
stored. Both the interaction training that follows and the content-baseline of
Section~\ref{sec:emb-baselines} read these precomputed vectors rather than
re-encoding the raw metadata, which keeps the product representations consistent
across experiments and avoids repeated encoding cost.

%--------------------------------------------------------------
\section{The Hierarchical Cosine-Softmax Objective}
\label{sec:emb-model}

\subsection{Intuition}
The core idea treats product interactions much like Word2Vec-style
skip-gram~\cite{Mikolov13} on item sequences: if users frequently move from
product $i$ to product $j$, then the embedding of $i$ should have high
similarity with the embedding of $j$. Rather than a flat softmax over the
entire catalogue, the probability of observing $j$ given $i$ is factored along
the hierarchy into three levels --- sub-category, then main-category, then
product --- each parameterised by cosine similarities between the source
embedding and the relevant category or product embeddings. This hierarchical
softmax~\cite{Morin05} makes the normalisation tractable at catalogue scale.

\subsection{Notation}
Let $e_p \in \mathbb{R}^{d}$ be the embedding of product $p$;
$c^{(\text{sub})}_s \in \mathbb{R}^{d}$ the embedding of sub-category $s$; and
$c^{(\text{mcat})}_m \in \mathbb{R}^{d}$ the embedding of main-category $m$.
Write $S$ for the set of all sub-categories, $M(s)$ for the main-categories
belonging to sub-category $s$, and $P_m$ for the products belonging to
main-category $m$. For a target product $j$, let $S_j$ and $M_j$ be its
sub-category and main-category. All embeddings are $L_2$-normalised on the
forward pass, so that $\lVert e_p\rVert_2 = \lVert c^{(\text{sub})}_s\rVert_2 =
\lVert c^{(\text{mcat})}_m\rVert_2 = 1$, and cosine similarity reduces to a dot
product, $\cos(x,y) = x\cdot y$.

\subsection{Probability Factorisation}
Given a source product $i$ and target product $j$, the three levels are:
\begin{align}
  P_{\text{sub}}(S_j \mid i)
    &= \frac{\exp\!\big(e_i^\top c^{(\text{sub})}_{S_j}\big)}
            {\sum_{s\in S}\exp\!\big(e_i^\top c^{(\text{sub})}_{s}\big)},
    \label{eq:emb-psub}\\[4pt]
  P_{\text{mcat}}(M_j \mid i, S_j)
    &= \frac{\exp\!\big(e_i^\top c^{(\text{mcat})}_{M_j}\big)}
            {\sum_{m\in M(S_j)}\exp\!\big(e_i^\top c^{(\text{mcat})}_{m}\big)},
    \label{eq:emb-pmcat}\\[4pt]
  P_{\text{prod}}(j \mid i, M_j)
    &= \frac{\exp\!\big(e_i^\top e_j\big)}
            {\sum_{k\in P_{M_j}}\exp\!\big(e_i^\top e_k\big)}.
    \label{eq:emb-pprod}
\end{align}
The main-category softmax is restricted to the main-categories under the target
sub-category $S_j$, and the product softmax is restricted to the products in
the target main-category $M_j$. The overall probability of observing $j$ given
$i$ is the product of the three:
\begin{equation}
  \label{eq:emb-ptotal}
  P_{\text{total}}(j\mid i)
    = P_{\text{sub}}(S_j\mid i)\,
      P_{\text{mcat}}(M_j\mid i,S_j)\,
      P_{\text{prod}}(j\mid i,M_j).
\end{equation}

\subsection{Loss and Dataset Objective}
For one observed interaction $(i\to j)$ the log-likelihood decomposes additively
across the three levels, and the per-sample loss is the negative log-likelihood
\begin{equation}
  \label{eq:emb-loss}
  \ell(i,j) = -\log P_{\text{total}}(j\mid i)
            = \ell_{\text{sub}}(i,j) + \ell_{\text{mcat}}(i,j)
              + \ell_{\text{prod}}(i,j),
\end{equation}
with the three terms
\begin{align}
  \ell_{\text{sub}}(i,j)  &= -e_i^\top c^{(\text{sub})}_{S_j}
     + \log\!\sum_{s\in S}\exp\!\big(e_i^\top c^{(\text{sub})}_{s}\big),
     \label{eq:emb-lsub}\\
  \ell_{\text{mcat}}(i,j) &= -e_i^\top c^{(\text{mcat})}_{M_j}
     + \log\!\sum_{m\in M(S_j)}\exp\!\big(e_i^\top c^{(\text{mcat})}_{m}\big),
     \label{eq:emb-lmcat}\\
  \ell_{\text{prod}}(i,j) &= -e_i^\top e_j
     + \log\!\sum_{k\in P_{M_j}}\exp\!\big(e_i^\top e_k\big).
     \label{eq:emb-lprod}
\end{align}
Let $\mathcal{D}_{\text{agg}}$ be the set of aggregated interaction pairs and
$w_{ij}$ the number of times transition $i\to j$ was observed after filtering.
The dataset-level objective is the weighted sum
\begin{equation}
  \label{eq:emb-objective}
  \mathcal{L}
    = \sum_{(i,j)\in\mathcal{D}_{\text{agg}}} w_{ij}\,\ell(i,j)
    = -\sum_{(i,j)\in\mathcal{D}_{\text{agg}}} w_{ij}\,
        \log P_{\text{total}}(j\mid i),
\end{equation}
minimised using SGD or Adam.

\subsection{Gradients}
Define the three softmax probabilities $p_{\text{sub}}(s\mid i)$,
$p_{\text{mcat}}(m\mid i,S_j)$, and $p_{\text{prod}}(k\mid i,M_j)$ exactly as in
\eqref{eq:emb-psub}--\eqref{eq:emb-pprod} but read as distributions over $s\in
S$, $m\in M(S_j)$, and $k\in P_{M_j}$ respectively. The gradient of the
per-sample loss with respect to the \emph{source-role} occurrence of $e_i$ is
\begin{equation}
  \label{eq:emb-grad-ei}
  \frac{\partial\ell(i,j)}{\partial e_i}
    = \Big(\sum_{s\in S} p_{\text{sub}}(s\mid i)\,c^{(\text{sub})}_{s}
           - c^{(\text{sub})}_{S_j}\Big)
    + \Big(\sum_{m\in M(S_j)} p_{\text{mcat}}(m\mid i,S_j)\,c^{(\text{mcat})}_{m}
           - c^{(\text{mcat})}_{M_j}\Big)
    + \Big(\sum_{k\in P_{M_j}} p_{\text{prod}}(k\mid i,M_j)\,e_k - e_j\Big),
\end{equation}
and for a \emph{target-role} occurrence $e_k$ with $k\in P_{M_j}$,
\begin{equation}
  \label{eq:emb-grad-ek}
  \frac{\partial\ell(i,j)}{\partial e_k}
    = \big(p_{\text{prod}}(k\mid i,M_j) - \mathbf{1}[k=j]\big)\,e_i.
\end{equation}
Because a single embedding table is used for both source and target roles, the
total gradient for a product embedding is the sum of all its contributions
across every interaction in which it appears as either a source or a target. In
the baseline the sub- and main-category centroids are fixed and frozen: they
participate in the forward pass but receive no gradient.

%--------------------------------------------------------------
\section{Extensions}
\label{sec:emb-extensions}

Two variants were developed beyond the frozen-centroid baseline.

\subsection{Trainable Centroids}
Here the loss form is unchanged, but $c^{(\text{sub})}_s$ and
$c^{(\text{mcat})}_m$ become learned parameters and therefore receive gradient
through their respective softmax denominators:
\begin{align}
  \frac{\partial\ell(i,j)}{\partial c^{(\text{sub})}_{s}}
    &= \big(p_{\text{sub}}(s\mid i) - \mathbf{1}[s=S_j]\big)\,e_i,
    \label{eq:emb-grad-csub}\\
  \frac{\partial\ell(i,j)}{\partial c^{(\text{mcat})}_{m}}
    &= \big(p_{\text{mcat}}(m\mid i,S_j) - \mathbf{1}[m=M_j]\big)\,e_i,
       \qquad m\in M(S_j).
    \label{eq:emb-grad-cmcat}
\end{align}
Main-category centroids outside $M(S_j)$ receive no gradient from the pair
$(i,j)$. Centroids are $L_2$-normalised on the forward pass exactly as products
are, though their master state is not constrained to unit norm between steps.

\subsection{Flat Objective}
The flat variant removes the hierarchical decomposition entirely, using a
single softmax over the whole product universe of size $N$:
\begin{equation}
  \label{eq:emb-flat}
  P(j\mid i) = \frac{\exp\!\big(e_i^\top e_j\big)}
                    {\sum_{k=1}^{N}\exp\!\big(e_i^\top e_k\big)},
  \qquad
  \ell(i,j) = -e_i^\top e_j
    + \log\!\sum_{k=1}^{N}\exp\!\big(e_i^\top e_k\big),
\end{equation}
with gradients, writing $p(k\mid i)=\exp(e_i^\top e_k)/\sum_t\exp(e_i^\top
e_t)$,
\begin{equation}
  \label{eq:emb-flat-grad}
  \frac{\partial\ell(i,j)}{\partial e_i} = \sum_{k=1}^{N} p(k\mid i)\,e_k - e_j,
  \qquad
  \frac{\partial\ell(i,j)}{\partial e_k}
    = \big(p(k\mid i) - \mathbf{1}[k=j]\big)\,e_i.
\end{equation}
The defining practical difference is the product-softmax denominator: the
hierarchical objective restricts it to $P_{M_j}$, so only products in the
target's main-category receive product-level gradient from a given pair,
whereas the flat objective places every product in every denominator, producing
a far denser gradient signal per step at correspondingly higher compute cost.

\subsection{Tied Embedding Table}
A single embedding table is shared between the source and target roles. For any
product that appears as both a source (via $e_i$) and a candidate (via $e_k$),
the total gradient is the sum of its source-role and candidate-role
contributions, applied once per optimisation step. This holds in both the
hierarchical and flat objectives, and is handled in the implementation by
merging the two gradient contributions before the Adam update.

%--------------------------------------------------------------
\section{Implementation}
\label{sec:emb-impl}

Training is organised into four phases. \textbf{Phase~A (interaction
preprocessing)} filters the interaction CSVs to valid product-to-product
transitions, maps them to embedding row indices, and aggregates them into
weighted pairs $(\text{src},\text{tgt},\text{count})$. \textbf{Phase~B
(taxonomy and centroids)} assigns each product to a sub- and main-category via
an admin taxonomy mapping and computes the sub- and main-category centroids as
the mean of member product embeddings. \textbf{Phase~C (embedding training)} is
the core optimisation and the locus of all experimental variation.
\textbf{Phase~D (recommendation generation)} $L_2$-normalises the final
embeddings, indexes them with a FAISS HNSW graph~\cite{FAISS,HNSW}, and queries
them to produce top-$K$ recommendations.

\paragraph{Optimiser and memory.}
All experiments use Adam with $\beta_1=0.9$, $\beta_2=0.999$,
$\epsilon=10^{-8}$, applied sparsely to the rows touched in each microbatch. The
embedding table and Adam moments are held in RAM as master state
($\approx 92$~GB for $20.2\text{M}\times 384$ across the parameter and the two
moment tensors), while per-training-block state is staged on the GPU, and disk
synchronisation occurs at the end of each epoch.

\paragraph{Performance rewrite.}
An initial memmap-backed implementation was bottlenecked by disk I/O and
host-device transfers, running at $\approx 2473$~s/epoch with the GPU at roughly
$15\%$ utilisation. Moving the master state to RAM, keeping per-block state
GPU-resident, and performing the Adam update on the GPU reduced this to
$\approx 780$~s/epoch --- a $3.2\times$ speed-up --- with no change to the
objective; training-loss curves before and after matched to within $0.1\%$ per
epoch, confirming the rewrite preserved the mathematics.

\paragraph{Correctness fixes.}
Two correctness-relevant fixes were made, both affecting the comparability of
results. First, \emph{sparse-update scope}: the optimised implementation
initially applied Adam densely to a block-level set of source rows each
microbatch, including rows with zero gradient; because Adam's update term is
nonzero even under zero gradient (the moment estimates decay at different
rates), this perturbed embeddings that should not have moved. It was fixed by
restricting each microbatch's Adam step to the rows that actually participated.
Second, \emph{tied-embedding coherence}: the first fix exposed a case where a
product appearing in both the source set and the candidate pool could have its
two GPU views diverge across microbatches; this was fixed by synchronising all
overlapping rows from the candidate-pool tensor back to the source tensor at the
end of each microbatch. Both fixes were in place for all trainable-centroid,
learning-rate-sweep, random-init, and industry-subset experiments; only the two
earliest full-catalog runs predate them and are approximately correct.

%--------------------------------------------------------------
\section{Experiments}
\label{sec:emb-experiments}

\subsection{Scale and Protocol}
The experiments were run at the scale summarised in
Table~\ref{tab:emb-scale}. All full-catalog runs share a catalogue of
$20{,}205{,}770\times 384$, Adam with $(\beta_1,\beta_2,\epsilon)=(0.9,0.999,
10^{-8})$, seed $42$, microbatch size $4096$, and a weighted-NLL loss.
Throughout, \emph{avg/pair} denotes the total weighted loss divided by the total
interaction weight, so that it is comparable across runs of differing size.

\begin{table}[ht]
  \centering
  \caption{Catalogue and data scale used in the embedding experiments
           (Phase~A aggregation).}
  \label{tab:emb-scale}
  \begin{tabular}{lr}
    \toprule
    Quantity & Value \\
    \midrule
    Products (embedding universe)     & $20{,}205{,}770$ \\
    Embedding dimension $d$           & $384$ \\
    Sub-categories                    & $687$ \\
    Main-categories                   & $98{,}353$ \\
    Aggregated interaction pairs      & $85{,}565{,}684$ \\
    Total interaction weight          & $130{,}702{,}326$ \\
    \bottomrule
  \end{tabular}
\end{table}

\subsection{Full-Catalog Results}
Table~\ref{tab:emb-full} collects the nine full-catalog runs: the two earliest
hierarchical runs (pre- and post-rewrite, frozen centroids), the two
trainable-centroid variants, a four-point product learning-rate sweep, and a
pure-random-initialisation run. The lowest total loss and the lowest avg/pair
of any full-catalog configuration are both achieved by the random-init run
(experiment~9).

\begin{table}[ht]
  \centering
  \caption{Full-catalog training results (weighted NLL). ``Cent.\ LR'' applies
           only to trainable-centroid runs. Best values in bold.}
  \label{tab:emb-full}
  \resizebox{\textwidth}{!}{%
  \begin{tabular}{clllccrr}
    \toprule
    \# & Variant & Init & Centroids & Prod.\ LR & Cent.\ LR & Epochs
       & Final total loss / avg/pair \\
    \midrule
    1 & Hierarchical (pre-rewrite) & text   & frozen    & $1\mathrm{e}{-3}$ & ---
      & 15 & $2.211\mathrm{e}{9}$ / $16.930$ \\
    2 & Frozen, optimized          & text   & frozen    & $1\mathrm{e}{-3}$ & ---
      & 10 & $2.2164\mathrm{e}{9}$ / $16.958$ \\
    3 & Trainable centroids v1     & text   & trainable & $1\mathrm{e}{-3}$
      & $1\mathrm{e}{-4}$ (both) & 10 & $2.1964\mathrm{e}{9}$ / $16.799$ \\
    4 & Trainable centroids v2     & text   & trainable & $1\mathrm{e}{-3}$
      & sub $1\mathrm{e}{-5}$, mcat $1\mathrm{e}{-4}$ & 10
      & $2.1975\mathrm{e}{9}$ / $16.813$ \\
    5 & LR sweep                   & text   & frozen    & $1\mathrm{e}{-3}$ & ---
      & 5  & $2.2305\mathrm{e}{9}$ / $17.065$ \\
    6 & LR sweep                   & text   & frozen    & $3\mathrm{e}{-3}$ & ---
      & 5  & $2.2177\mathrm{e}{9}$ / $16.968$ \\
    7 & LR sweep                   & text   & frozen    & $1\mathrm{e}{-2}$ & ---
      & 5  & $2.2118\mathrm{e}{9}$ / $16.922$ \\
    8 & LR sweep                   & text   & frozen    & $5\mathrm{e}{-2}$ & ---
      & 5  & $2.2150\mathrm{e}{9}$ / $16.947$ \\
    9 & Random init                & random & trainable & $1\mathrm{e}{-2}$
      & sub $1\mathrm{e}{-4}$, mcat $1\mathrm{e}{-3}$ & 10
      & $\mathbf{2.1831\mathrm{e}{9}}$ / $\mathbf{16.703}$ \\
    \bottomrule
  \end{tabular}%
  }
\end{table}

\subsection{Component-Split Diagnostic}
A key diagnostic logged the three loss terms separately.
Table~\ref{tab:emb-split} gives representative values for the frozen-centroid,
text-init run (experiment~2). Across every full-catalog configuration,
$\ell_{\text{sub}}$ and $\ell_{\text{mcat}}$ are nearly flat --- indeed
$\ell_{\text{sub}}$ tends to drift slightly \emph{upward} --- while
$\ell_{\text{prod}}$ accounts for essentially all of the loss reduction.

\begin{table}[ht]
  \centering
  \caption{Per-component loss over epochs (experiment~2: text init, frozen
           centroids).}
  \label{tab:emb-split}
  \begin{tabular}{cccc}
    \toprule
    Epoch & $\ell_{\text{sub}}$ & $\ell_{\text{mcat}}$ & $\ell_{\text{prod}}$ \\
    \midrule
    1  & $8.041\mathrm{e}{8}$ & $6.652\mathrm{e}{8}$ & $8.121\mathrm{e}{8}$ \\
    5  & $7.982\mathrm{e}{8}$ & $6.637\mathrm{e}{8}$ & $7.685\mathrm{e}{8}$ \\
    10 & $8.015\mathrm{e}{8}$ & $6.644\mathrm{e}{8}$ & $7.505\mathrm{e}{8}$ \\
    \bottomrule
  \end{tabular}
\end{table}

\subsection{Random Initialisation}
In this run the embeddings were replaced with $\mathcal{N}(0,\sigma^2)$ samples,
$\sigma=0.02$, drawn deterministically (seed $42$) at the same
$20{,}205{,}770\times 384$ shape, with centroids recomputed from the random
embeddings. Three observations stand out: the random-init epoch-1
$\ell_{\text{prod}}$ ($7.80\mathrm{e}{8}$) was already \emph{lower} than the
text-init epoch-1 $\ell_{\text{prod}}$ ($8.12\mathrm{e}{8}$);
$\ell_{\text{prod}}$ reached its minimum by epoch~4 and then crept upward
slightly; and the run produced the lowest total loss of any full-catalog
configuration.

\subsection{Industry-Subset Experiments (Furniture)}
To enable faster iteration and to test the architecture at smaller scale, the
pipeline was restricted to a single industry (Furniture) by filtering and
reindexing while preserving the original taxonomy --- $608{,}287$ products,
$2{,}065{,}519$ aggregated pairs, and $1{,}305$ main-categories.
Table~\ref{tab:emb-furniture} contrasts the hierarchical objective with the
flat objective under identical data, initialisation, and learning rate.

\begin{table}[ht]
  \centering
  \caption{Furniture-subset training results: hierarchical versus flat
           objective.}
  \label{tab:emb-furniture}
  \begin{tabular}{cllcr}
    \toprule
    \# & Variant & Objective & Epochs & Final avg/pair \\
    \midrule
    10 & Furniture hierarchical & 3-term hierarchical     & 50 & $13.204$ \\
    11 & Furniture flat         & flat over $608{,}287$    & 30 & $12.354$ \\
    \bottomrule
  \end{tabular}
\end{table}

The gap is stark. The hierarchical subset run moved avg/pair by only $0.030$
over $50$ epochs ($13.234\to 13.204$), whereas the flat run moved it by $0.300$
over $30$ epochs ($12.654\to 12.354$) --- with the epoch-1 flat loss already
below the hierarchical run's converged value. At convergence the flat model
sits roughly one nat below the uniform-random baseline of
$\log(608{,}287)\approx 13.32$.

%--------------------------------------------------------------
\section{Findings}
\label{sec:emb-findings}

The training-loss evidence points to three consistent conclusions. Each
concerns optimisation behaviour only; recommendation-quality evaluation is
required to confirm whether any of them translates into better retrieval.

\begin{enumerate}
  \item \textbf{The hierarchical structure contributes little learning signal.}
        Across all full-catalog configurations the sub- and main-category loss
        terms are effectively static while the product term does nearly all the
        work (Table~\ref{tab:emb-split}). The three-level decomposition behaves
        as a denominator-sharing \emph{computational} structure rather than as a
        source of useful gradient. Frozen centroids do not by themselves imply
        static sub/mcat loss: those terms still move because the source
        embeddings $e_i$ scored against the fixed centroids are themselves
        changing each step --- in the text-init case, drifting away from their
        sub-category centroid and slowly raising $\ell_{\text{sub}}$.
  \item \textbf{Text-based initialisation is a head start, not an optimum.}
        Random initialisation reached a lower total loss than text
        initialisation, and did so with a lower product loss even in the first
        epoch. The geometry imposed by product text is therefore correlated
        with, but distinct from, the geometry that best predicts interactions;
        training spends effort relaxing the text geometry toward the interaction
        geometry.
  \item \textbf{Flat softmax materially outperforms the hierarchical objective
        at subset scale.} Under identical data, initialisation, and learning
        rate, the flat objective achieved roughly an order of magnitude more
        loss reduction than the hierarchical objective on the Furniture subset.
        The hierarchical decomposition, attractive for tractability at full
        catalogue scale, appears to \emph{impede} convergence when the universe
        is small enough for a full softmax to be computed directly.
\end{enumerate}

All three are statements about the training objective. Two of them --- the
lower loss of random initialisation and the superiority of the flat objective
--- turn out \emph{not} to carry over to recommendation quality, as the
click-through evaluation of Section~\ref{sec:emb-ctr} shows.

For reference, Table~\ref{tab:emb-hparams} lists the hyperparameters used across
the experiments. All runs are seeded ($42$) for embedding initialisation and
microbatch shuffling; random embeddings are generated deterministically from
\texttt{numpy.random.default\_rng(42)} at the catalogue shape and frozen to
disk; master state (embeddings and Adam moments) is synchronised to disk at each
epoch end for exact resume; and per-epoch summaries are written to a
\texttt{train\_state.json} for every run.

\begin{table}[ht]
  \centering
  \caption{Hyperparameter reference for the embedding experiments.}
  \label{tab:emb-hparams}
  \begin{tabular}{ll}
    \toprule
    Hyperparameter & Value(s) used \\
    \midrule
    Embedding dimension $d$          & $384$ \\
    Optimiser                        & Adam \\
    $\beta_1,\beta_2,\epsilon$       & $0.9,\ 0.999,\ 10^{-8}$ \\
    Product learning rate            & $1\mathrm{e}{-3},\ 3\mathrm{e}{-3},\ 1\mathrm{e}{-2},\ 5\mathrm{e}{-2}$ \\
    Sub-category centroid LR         & $1\mathrm{e}{-4}$ (v1), $1\mathrm{e}{-5}$ (v2), $1\mathrm{e}{-4}$ (random init) \\
    Main-category centroid LR        & $1\mathrm{e}{-4}$ (v1), $1\mathrm{e}{-4}$ (v2), $1\mathrm{e}{-3}$ (random init) \\
    Microbatch size (hierarchical)   & $4096$ pairs per main-category block \\
    Microbatch size (flat)           & $1024$ pairs \\
    Epochs                           & $5\,/\,10\,/\,15\,/\,30\,/\,50$ (per experiment) \\
    Random seed                      & $42$ \\
    Random-init distribution         & $\mathcal{N}(0,\,0.02^2)$ \\
    Weight scale                     & $1.0$ \\
    \bottomrule
  \end{tabular}
\end{table}

%--------------------------------------------------------------
\section{Recommendation-Quality Evaluation}
\label{sec:emb-ctr}

The findings above concern the training objective. We now measure what the
embeddings actually deliver as recommendations, using offline click-through
rate (CTR) on the four held-out weeks of June~2025 (training window
1~July~2024--31~May~2025). CTR is reported on two surfaces: \emph{IMOB}, the
IndiaMART mobile app, and \emph{IM}, the IndiaMART web surface. At serving time
the learned embeddings are $L_2$-normalised, indexed with FAISS HNSW, and
queried for top-$K$ neighbours, after which a geographic filter is applied
(retaining products within roughly $50$~km for Tier-1 and $250$~km for Tier-2
cities).

\subsection{Baselines}
\label{sec:emb-baselines}
Three baselines are used. The \emph{Frequency-Based Model} (FBM), already used as
the production baseline in Section~\ref{sec:mf-ctr-setup}, is the incumbent
frequency/co-occurrence recommender: an $n\times n$ matrix over products in which
entry $m_{ij}$ is the number of observed transitions from product $i$ to product
$j$, with each row normalised to sum to one. \emph{Backfilling}
(Section~\ref{sec:mf-ctr-setup}) is a coverage mechanism layered on top --- for a product with
fewer than $k$ recommendations, or with recommendation scores below a
threshold, the remaining slots are filled with IndiaMART's pre-existing
recommendations. FBM is therefore reported both \emph{without} and \emph{with}
backfilling, the latter being the strongest baseline. The third baseline, the
\emph{Content-Baseline Model}, is the text-metadata cold-start embeddings served
directly, with no interaction training --- it isolates how much the skip-gram
training over interactions adds on top of the initial text geometry.

\subsection{Main Results}
\label{sec:emb-main-ctr}
Table~\ref{tab:emb-ctr-main} compares the baselines against the hierarchical
cosine-softmax model (frozen centroids, text initialisation) at epochs~1, 5,
and~10, on the mobile surface.

\begin{table}[ht]
  \centering
  \caption{Offline CTR (\%, IMOB / mobile) on held-out weeks of June~2025:
           baselines versus the frozen hierarchical softmax model.}
  \label{tab:emb-ctr-main}
  \small
  \begin{tabular}{lcccc}
    \toprule
    Model & 01~Jun & 08~Jun & 15~Jun & 22~Jun \\
    \midrule
    FBM without backfilling      & $8.12$ & $7.52$ & $7.00$ & $6.58$ \\
    Content-Baseline (text only) & $4.07$ & $3.83$ & $3.81$ & $3.80$ \\
    FBM with backfilling         & $\mathbf{10.83}$ & $\mathbf{10.22}$ & $\mathbf{9.96}$ & $\mathbf{9.54}$ \\
    \midrule
    Hierarchical softmax, epoch~1  & $5.93$ & $5.67$ & $5.48$ & $5.33$ \\
    Hierarchical softmax, epoch~5  & $7.22$ & $6.94$ & $6.61$ & $6.34$ \\
    Hierarchical softmax, epoch~10 & $7.24$ & $6.96$ & $6.61$ & $6.34$ \\
    \bottomrule
  \end{tabular}
\end{table}

Three observations follow. First, interaction training clearly helps: CTR rises
from $5.93\%$ at epoch~1 to $7.24\%$ at epoch~10 (week of 1~June), and the
learned model comfortably beats the content-only baseline ($\approx 4\%$),
confirming that the skip-gram objective adds real signal over the text geometry.
Second, the gains saturate quickly --- epochs~5 and~10 are almost
indistinguishable --- consistent with the training-loss curves. Third, and less
comfortably, the embedding model does \emph{not} overtake the frequency
baseline: it trails FBM without backfilling in every week and sits well below
FBM with backfilling (the $\approx 9.5$--$10.8\%$ incumbent). Learned
embeddings, on their own, do not yet match the production frequency model.

\subsection{Centroid Variants and Initialisation}
\label{sec:emb-variants-ctr}
Tables~\ref{tab:emb-ctr-imob} and~\ref{tab:emb-ctr-im} report the trainable-%
centroid variants and the random-initialisation run on the two surfaces.
Trainable centroids (v1, v2) track the frozen model almost exactly --- the best
single cell is v1 at epoch~5 ($7.36\%$), a marginal edge --- reinforcing the
training-loss finding that the centroid machinery adds little. The
random-initialisation run is the telling case: although it achieved the
\emph{lowest training loss} of any full-catalog configuration
(Section~\ref{sec:emb-experiments}), its CTR is markedly \emph{worse} than
text-initialisation, at $6.18\%$ (epoch~10, mobile, week of 1~June) against the
frozen model's $7.24\%$, and far worse at epoch~1 ($3.48\%$ versus $5.93\%$).
Lower loss, worse recommendations.

\begin{table}[ht]
  \centering
  \caption{Offline CTR (\%, IMOB / mobile) for the full-catalog embedding
           variants on held-out weeks of June~2025.}
  \label{tab:emb-ctr-imob}
  \small
  \begin{tabular}{lcccc}
    \toprule
    Variant (epoch) & 01~Jun & 08~Jun & 15~Jun & 22~Jun \\
    \midrule
    Frozen centroids, e1   & $5.93$ & $5.67$ & $5.48$ & $5.33$ \\
    Frozen centroids, e5   & $7.22$ & $6.94$ & $6.61$ & $6.34$ \\
    Frozen centroids, e10  & $7.24$ & $6.96$ & $6.61$ & $6.34$ \\
    Trainable v1, e1       & $6.25$ & $5.98$ & $5.74$ & $5.58$ \\
    Trainable v1, e5       & $7.36$ & $7.03$ & $6.76$ & $6.53$ \\
    Trainable v1, e10      & $7.20$ & $6.87$ & $6.59$ & $6.35$ \\
    Trainable v2, e5       & $6.03$ & $5.76$ & $5.55$ & $5.37$ \\
    Trainable v2, e10      & $7.26$ & $6.93$ & $6.64$ & $6.40$ \\
    Random init, e1        & $3.48$ & $3.35$ & $3.06$ & $2.85$ \\
    Random init, e5        & $6.05$ & $5.81$ & $5.50$ & $5.20$ \\
    Random init, e10       & $6.18$ & $5.93$ & $5.65$ & $5.35$ \\
    \bottomrule
  \end{tabular}
\end{table}

\begin{table}[ht]
  \centering
  \caption{Offline CTR (\%, IM / web) for the full-catalog embedding variants on
           held-out weeks of June~2025. (Frozen-model web CTR was reported only
           under the geographic slices of Table~\ref{tab:emb-ctr-delhi}.)}
  \label{tab:emb-ctr-im}
  \small
  \begin{tabular}{lcccc}
    \toprule
    Variant (epoch) & 01~Jun & 08~Jun & 15~Jun & 22~Jun \\
    \midrule
    Trainable v1, e1   & $5.37$ & $5.13$ & $5.00$ & $4.91$ \\
    Trainable v1, e5   & $5.87$ & $5.42$ & $5.38$ & $5.46$ \\
    Trainable v1, e10  & $5.45$ & $5.05$ & $4.96$ & $4.91$ \\
    Trainable v2, e5   & $5.23$ & $4.95$ & $4.89$ & $4.94$ \\
    Trainable v2, e10  & $5.52$ & $5.09$ & $5.00$ & $4.93$ \\
    Random init, e1    & $1.73$ & $1.65$ & $1.67$ & $1.93$ \\
    Random init, e5    & $3.77$ & $3.45$ & $3.49$ & $3.85$ \\
    Random init, e10   & $4.00$ & $3.70$ & $3.69$ & $3.77$ \\
    \bottomrule
  \end{tabular}
\end{table}

\subsection{Geographic Slices}
\label{sec:emb-geo-ctr}
Because recommendations are geographically filtered at serving time, CTR was
also computed under two Delhi-scoped evaluation filters, denoted \emph{Delhi
strict} and \emph{Delhi cluster}, for the frozen hierarchical model.
Table~\ref{tab:emb-ctr-delhi} gives both, on both surfaces, at epochs~1, 5,
and~10. The absolute levels are lower than the unfiltered headline (the filters
restrict the candidate pool), but the qualitative picture is unchanged: CTR
improves with training and mobile leads web.

\begin{table}[ht]
  \centering
  \caption{Offline CTR (\%) of the frozen hierarchical softmax model under Delhi
           geographic filters, by surface and epoch, June~2025.}
  \label{tab:emb-ctr-delhi}
  \small
  \resizebox{\textwidth}{!}{%
  \begin{tabular}{llccccccccc}
    \toprule
    & & \multicolumn{3}{c}{Delhi strict} & & \multicolumn{3}{c}{Delhi cluster} \\
    \cmidrule(lr){3-5}\cmidrule(lr){7-9}
    Surface & Week & e1 & e5 & e10 & & e1 & e5 & e10 \\
    \midrule
    IMOB & 01~Jun & $6.15$ & $6.36$ & $5.96$ & & $5.44$ & $5.83$ & $5.60$ \\
    IMOB & 08~Jun & $5.93$ & $6.51$ & $6.24$ & & $5.46$ & $5.90$ & $5.71$ \\
    IMOB & 15~Jun & $6.11$ & $6.45$ & $6.18$ & & $5.52$ & $5.79$ & $5.59$ \\
    IMOB & 22~Jun & $6.40$ & $6.80$ & $6.54$ & & $5.56$ & $5.83$ & $5.56$ \\
    \midrule
    IM   & 01~Jun & $5.65$ & $5.37$ & $5.08$ & & $5.37$ & $5.44$ & $4.82$ \\
    IM   & 08~Jun & $5.28$ & $5.05$ & $4.85$ & & $4.29$ & $4.08$ & $3.80$ \\
    IM   & 15~Jun & $5.75$ & $5.51$ & $5.09$ & & $4.59$ & $4.54$ & $4.27$ \\
    IM   & 22~Jun & $5.75$ & $4.83$ & $4.39$ & & $4.94$ & $5.33$ & $5.06$ \\
    \bottomrule
  \end{tabular}%
  }
\end{table}

\subsection{Industry Subset: Flat versus Hierarchical}
\label{sec:emb-industry-ctr}
The most instructive comparison is on the Furniture subset, where the flat and
hierarchical objectives were trained under identical conditions.
Tables~\ref{tab:emb-ctr-flat} and~\ref{tab:emb-ctr-hier} report CTR for the two.
Recall that on the training objective the flat model was the clear winner
(Section~\ref{sec:emb-findings}). On recommendation quality the ranking
\emph{inverts}. The hierarchical model climbs from $3.78\%$ (epoch~1, mobile,
week of 1~June) to $7.10\%$ by epoch~50, whereas the flat model peaks at
$4.91\%$ at epoch~1 and then \emph{degrades} with further training, settling
around $3.0\%$ --- worse than where it started. The configuration with the lower
training loss produces the worse, and worsening, recommendations.

\begin{table}[ht]
  \centering
  \caption{Furniture-subset offline CTR (\%) for the \emph{flat} objective, by
           surface and epoch, June~2025.}
  \label{tab:emb-ctr-flat}
  \small
  \begin{tabular}{llccccc}
    \toprule
    Surface & Week & e1 & e5 & e10 & e20 & e30 \\
    \midrule
    IMOB & 01~Jun & $4.91$ & $3.07$ & $2.98$ & $3.00$ & $3.00$ \\
    IMOB & 08~Jun & $4.76$ & $2.95$ & $2.94$ & $2.91$ & $2.93$ \\
    IMOB & 15~Jun & $4.65$ & $2.86$ & $2.81$ & $2.87$ & $2.89$ \\
    IMOB & 22~Jun & $4.55$ & $2.77$ & $2.76$ & $2.82$ & $2.82$ \\
    \midrule
    IM   & 01~Jun & $3.68$ & $2.17$ & $2.19$ & $2.35$ & $2.38$ \\
    IM   & 08~Jun & $3.13$ & $2.03$ & $1.92$ & $1.81$ & $1.73$ \\
    IM   & 15~Jun & $3.15$ & $1.96$ & $1.88$ & $1.88$ & $1.96$ \\
    IM   & 22~Jun & $3.35$ & $1.96$ & $2.16$ & $2.15$ & $2.05$ \\
    \bottomrule
  \end{tabular}
\end{table}

\begin{table}[ht]
  \centering
  \caption{Furniture-subset offline CTR (\%) for the \emph{hierarchical}
           objective, by surface and epoch, June~2025.}
  \label{tab:emb-ctr-hier}
  \small
  \begin{tabular}{llcccc}
    \toprule
    Surface & Week & e1 & e10 & e25 & e50 \\
    \midrule
    IMOB & 01~Jun & $3.78$ & $6.83$ & $7.13$ & $7.10$ \\
    IMOB & 08~Jun & $3.62$ & $6.54$ & $6.86$ & $6.87$ \\
    IMOB & 15~Jun & $3.54$ & $6.28$ & $6.58$ & $6.61$ \\
    IMOB & 22~Jun & $3.42$ & $6.11$ & $6.33$ & $6.43$ \\
    \midrule
    IM   & 01~Jun & $2.80$ & $4.13$ & $4.29$ & $4.29$ \\
    IM   & 08~Jun & $2.56$ & $3.66$ & $3.89$ & $3.92$ \\
    IM   & 15~Jun & $2.46$ & $3.67$ & $3.75$ & $3.77$ \\
    IM   & 22~Jun & $2.65$ & $3.91$ & $4.00$ & $3.98$ \\
    \bottomrule
  \end{tabular}
\end{table}

\subsection{Discussion: Training Loss versus Recommendation Quality}
\label{sec:emb-discussion}
The evaluation delivers a consistent and cautionary message. In the two places
where the training objective and recommendation quality could be compared
head-to-head, they disagreed. Random initialisation reached a lower weighted NLL
than text initialisation yet produced lower CTR; the flat objective reached a
much lower NLL than the hierarchical objective on the Furniture subset yet
produced substantially lower CTR, degrading as its loss improved. Weighted
negative log-likelihood is thus, in this setting, an unreliable proxy for
recommendation quality --- driving it down can leave retrieval unchanged or make
it worse.

Two practical conclusions follow. First, text-based initialisation is valuable
for \emph{recommendations} even though it is not the loss optimum: the geometry
inherited from product text places items sensibly from the outset, and training
away from it (as random init does) trades a better loss for worse retrieval.
Second, the hierarchical decomposition, though it contributes little gradient
signal and is beaten on loss by the flat softmax, yields better recommendations
at subset scale --- its restriction of the product softmax to within-category
candidates appears to act as a useful inductive constraint rather than merely a
computational convenience. Overall, the best embedding configuration (text-%
initialised hierarchical softmax, $\approx 7.2\%$ CTR on mobile) is a clear
improvement over the content-only baseline but remains below the
frequency-based incumbent with backfilling, which motivates the further
directions explored elsewhere in this thesis.
